\documentclass[output=paper,colorlinks,citecolor=brown]{langscibook} 

\bibliography{localbibliography}

%This is where you put the authors and their affiliations
\author{Your Name Here\affiliation{Pomona College}
}

%Insert your title here
\title{Field Methods Paper Template}  
\abstract{Here is where you write the abstract of the paper.
}

\input{FieldMethodsTemplatePackages.tex}
\input{FieldMethodsTemplateCommands.tex}

% \IfFileExists{../localcommands.tex}{%hack to check whether this is being compiled as part of a collection or standalone
%   \input{../localpackages}
%   \input{../localcommands} 
% \togglepaper[23]
% }{}

\begin{document}
\maketitle

%\tableofcontents 

\section{What is this document?}

%Please note, this template is built on Language Science Press's LaTeX template from Overleaf (as of early 2020). It has been posted with their permission, but any errors should be attributed to Diercks and nobody else. If you intend to write for Language Science Press independently, please download their LaTeX templates directly from their website. http://langsci-press.org/

\begin{itemize}

\item This is a Paper template for writing papers in \LaTeX{} for Field Methods classes, distributed as \href{https://www.overleaf.com/latex/templates/field-methods-paper-template/pqrxmrsptqcv}{an Overleaf template}.

\item If you want a more in-depth introduction to \LaTeX{}, see the \href{https://www.overleaf.com/latex/templates/field-methods-papers-quick-reference-guide/cqbjxwnffvrm}{Field Methods Quick Reference Guide}.

\item If you clicked on ``\href{https://www.overleaf.com/latex/templates/field-methods-paper-template/pqrxmrsptqcv}{Open as a Template},'' a copy of this project is now in your Overleaf account (i.e. it belongs to you), and you can now edit this document to be an actual paper that you write.\footnote{This also means you will probably want to rename the document - you can do this by hovering your cursor over the title ``Field Methods Paper Template' at the top of this page, and then clicking on the pencil icon that appears to the right of the title. You can also rename the FieldMethodsTemplate.tex file in the leftmost panel by clicking on the down arrow.} 

\item Designing a \LaTeX{} project from scratch can be difficult; the point of this template is that this is already done for you, you can just start writing. 

\item Don't worry about breaking anything - you can always open the template link again from Overleaf's website to see this initial version again. So edit away!


\end{itemize}

\section{Basic commands, a reminder}

\subsection{Text formatting commands}

\begin{table}
\caption{Selected commands for common ACAL formatting needs}
\label{tab:diercks:CommandsTable}

\begin{tabular}[t]{l l l}
\lsptoprule		
\textbf{Symbol/Annotation} & \textbf{Example} & \textbf{Code} \\ \midrule
		Ellipsis & \dots & \verb|\dots| \\
		Subscript & NP\subs{i} & \verb|NP\subs{i}| \\
		Superscript & NP\supers{i} & \verb|NP\supers{i}|\\ 
		Bold & \textbf{bold} & \verb|\textbf{bold}| \\
		Italic & \textit{italic} & \verb|\textit{italic}| \\
		Small Caps & \textsc{small caps} & \verb|\textsc{small caps}| \\
		Strikeout & \sout{strikeout} & \verb|\sout{strikeout}| \\
		Underline & \uline{underline} & \verb|\uline{underline}| \\
		circle something in text & \circled{something} & \verb|\circled{something}| \\
		Highlight something & \hl{something} & \verb|\hl{something}| \\
		Null & \nothing  & \verb|\nothing|\\
	\lspbottomrule
	\end{tabular}
\end{table}

\subsection{Numbered examples} \label{sec:diercks:NumberedExs}

\ea \label{ex:diercks:BukusuNegation}
%\langinfo{Lubukusu}{}{Justine Sikuku, pc}
\gll Wekesa se-a-la-ba a-kula ka-ma-indi ta. \\
Wekesa \Neg-1\Sm-\Fut-be 1\Sm-buy 6-6-maize \Neg{} \\ 
\glt `Wekesa will not be buying maize.'
\z

\ea \label{ex:diercks:XlistExample}
\begin{xlist}

\ex This is the first example. 

\ex \label{ex:diercks:SwahiliXlist}
\gll hii ni m-fano w-a pili, kwa Ki-swahili \\
this is 3-example 3-\Assoc{} second, of 7-Swahili \\
\glt `This is a second example, in Swahili.' 

\end{xlist}

\z 

\section{Using cross-references: Different from standard \LaTeX{}} \label{sec:diercks:CrossReferences}

\begin{enumerate}    

\item These are Language Science Press's commands for creating labels (different for what sort of element you are labeling): 

\begin{itemize}

\item Section label: \verb|\label{sec:author:yourlabel}|

\item Table label: \verb|\label{tab:author:yourlabel}|

\item Figure label: \verb|\label{fig:author:yourlabel}|

\item Example label: \verb|\label{ex:author:yourlabel}|

\end{itemize}
    
    
\item These are LSP's commands for creating cross-references (again, different for what kind of element is being referenced). 

\begin{itemize}

\item Section reference (will read `\S \#’ in output):\\ \verb|\sectref{sec:author:yourlabel}|

\item Table reference (will read `Table \#’ in output):\\ \verb|\tabref{tab:author:yourlabel}|

\item Figure reference (will read `Figure \#’ in output): \\ \verb|\figref{fig:author:yourlabel}|

\item Example reference with parens — note caps (will read `(\#)’ in output): \\ \verb|\REF{ex:author:yourlabel}| 

\item Example reference without parens — note caps (will read `\#’ in output): \\ \verb|\ref{ex:author:yourlabel}|

\end{itemize}

\item For examples, see the \href{https://www.overleaf.com/latex/templates/field-methods-papers-quick-reference-guide/cqbjxwnffvrm}{Field Methods Quick Reference  Guide}.

\end{enumerate}


\section{LSP Table requirements}

LSP has a few specific restrictions on how to structure tables: 

\begin{enumerate}
    
    \item No vertical lines in tables
    
    \item Tables (created via tabular) are created as \LaTeX{} tables (i.e. using
    \verb|\begin{table}|, and are \textit{not} created as numbered examples, though this is otherwise common in linguistics.
    
    \item There is an LSP-specific command for the toprule (\verb|\lsptoprule|) and bottomrule (\verb|\lspbottomrule|) of a table, exemplified in the table below

\end{enumerate}
 
\begin{table}

  \caption{A sample table}
  \label{tab:diercks:SampleTableHeaders}
  
\begin{tabular}{lll}
\lsptoprule %This (\lsptoprule) is a required piece of formatting for tables in LSP volumes 
Header 1 & Header 2 & Header 3 \\
\hline\hline
cell 1 & cell 2 & cell 3 \\
\hline
cell 4 & cell 5 & cell 6 \\
\hline
cell 7 & cell 8 & cell 9 \\
\lspbottomrule %This (\lspbottomrule) is a required piece of formatting for tables in LSP volumes 


\end{tabular}
\end{table}

\section*{Abbreviations}

To be added later.

%Abbreviations should be included in this section. It is a required part of volumes with Language Science Press that an abbreviations list be included, so please do not remove this section, but for Field Methods you don't have to write anything here.

\section*{Acknowledgements}

Acknowledgements should be listed here.

\noindent This \href{https://www.overleaf.com/latex/templates/field-methods-paper-template/pqrxmrsptqcv}{Paper Template} and the \href{https://www.overleaf.com/latex/templates/field-methods-papers-quick-reference-guide/cqbjxwnffvrm}{corresponding Quick Reference Guide} borrow directly (with permission) from corresponding documents written by Michael Diercks and Franny Brogan for linguistics students at Pomona College, and from the in-development \LaTeX{} template for proceedings papers for ACAL (Annual Conference on African Linguistics). These templates are built on Language Science Press's Overleaf template (for eventual publishing with LSP), and are published on Overleaf with LSP permission. Our thanks to Sebastian Nordhoff for his feedback on earlier versions of this template. Diercks' original \LaTeX{} teaching materials were created with assistance from Claire Halpert, Nico Baier, Jason Zentz, Maddy Bossi, and Michael Clausen.

\printbibliography[heading=subbibliography,notkeyword=this]

\end{document}
